\begin{abstractpage}

\begin{abstract}{romanian}
Acestă lucrare prezintă implementarea unei platforme web interactive destinată învățării și practicării programării competitive. Tema este relevantă deoarece folosește instrumente moderne de e-learning care să combată lipsa de motivare folosind tehnici de gamificare și interacțiune socială. Din punct de vedere tehnologic, lucrarea utilizează o arhitectură client-server bazată pe framework-ul Spring Boot pentru backend și o bază de date relațională gestionată prin Hibernate. Frontend-ul este scris în React cu TypeScript utilizând TanStack Query pentru administrarea eficientă a stărilor și Material UI pentru interfață. Rezultate obținute constă într-o platformă funcțională care integreaza un motor de execuție automată a codului, un magazin de efecte vizuale și un modul social pentru interacțiunea utilizatorilor.
\end{abstract}

\begin{abstract}{english}
This thesis presents the implementation of a web application used for learning and practising competitive programming. This topic is relevant because it uses modern e-learning methods to prevent the loss of motivation by using gamification and social interactions techniques. Technologically, the thesis uses a client-server architecture based on the Spring Boot framework on backend and a relational database managed using Hibernate. The frontend has been developed using React with TypeScript and TanStack Query for the efficient management of states and Material UI for interface. The obtained results consist of a fully functional platform that integrates an automated code execution engine, a shop for visual effects and a social module for user interaction. 
\end{abstract}

\end{abstractpage}